\documentclass[11pt]{article}

\usepackage[margin=1in]{geometry}
\usepackage{amsmath,amssymb,amsthm,mathtools}
\usepackage{booktabs}
\usepackage{enumitem}
\usepackage{microtype}
\usepackage[numbers,sort&compress]{natbib}
\usepackage[hidelinks]{hyperref}
\usepackage{xcolor}

\newtheorem{definition}{Definition}[section]
\newtheorem{theorem}[definition]{Theorem}
\newtheorem{proposition}[definition]{Proposition}
\newtheorem{corollary}[definition]{Corollary}
\newtheorem{principle}[definition]{Design Principle}
\newtheorem{conjecture}[definition]{Conjecture}
\newtheorem{remark}[definition]{Remark}

\newcommand{\Reach}{\operatorname{Reach}}
\newcommand{\Thm}{\operatorname{Thm}}
\newcommand{\Verify}{\mathsf{Verify}}
\newcommand{\E}{\mathbb{E}}
\newcommand{\Prob}{\mathbb{P}}
\newcommand{\ind}{\mathbf{1}}
\newcommand{\bits}{\operatorname{bits}}

\title{Guided Proof-Space Evolution:\\
Finite Proof Radius and Completeness-Preserving Resource Allocation\\
for AI-Assisted Mathematics}

\ifdefined\anonymousmode
  \author{Anonymous Author}
\else
  \author{Alex Chengyu Li\\
  \small GenAI Architect; Independent Researcher, Tokyo, Japan\\
  \small \texttt{cl7076@nyu.edu}}
\fi
\date{August 2026}

\begin{document}
\maketitle

\begin{abstract}
A theorem provable in a declared machine-checkable system has a finite proof
code.  Exhaustive search therefore has a finite target radius on every positive
instance.  Yet when the theorem set is undecidable, no computable uniform bound
reveals that radius in advance.  This tension creates the operational problem
for AI-assisted mathematics: within a fixed proof system, guidance cannot
change which proof objects are valid, but it can change by exponential factors
how much probability and compute reach them.

This paper develops a proof-search control plane for that problem.  A complete
proposal protocol induces probability mass on candidate proof objects.  We
define a target's policy-relative distance as the negative logarithm of the
total mass assigned to its valid proofs.  Under independent verifier-guided
sampling, expected candidate count is the inverse of that proof mass.  A unit
reduction in distance therefore yields a multiplicative search gain.  Human
judgment, language-model guidance, route selection, and tool choice enter as
probability redistribution, not as evidentiary authority.

To prevent guidance from deleting the only valid route, we mix an arbitrary
adaptive policy with a support-complete search process.  This
\emph{completeness reserve} preserves almost-sure discovery of every finite
proof and gives an explicit overhead bound; we also prove a deterministic
dovetailing analogue.  We then lift the model from one target to a portfolio of
typed claims and heterogeneous tools.  The scheduler values an action by its
probability of earning a proof edge, its downstream reachability gain, and its
generation, verification, transport, and composition costs.  For a homogeneous
conjunctive route, an exact one-unit allocation law shows that moving one
attempt from an edge with at least two more attempts to the less-funded edge
strictly increases closure probability.  Candidate artifacts change no
mathematical reachability until adequate checking and exact claim binding make
their edges admissible.

The resulting control plane is completeness-preserving and
verification-cost-aware.  Its central object is not an isolated proof tree but
the policy that converts heterogeneous search into earned, portfolio-level
reachability.  The framework composes universal search, neural guidance,
claim-bound checking, and proof engineering while preserving their distinct
guarantees.
\end{abstract}

\section{Introduction}

AI systems can now generate proof sketches, tactic sequences, programs,
conjectures, counterexamples, and formal fragments at a rate that changes the
organization of mathematical work.  Formal theorem provers increasingly pair
a learned proposal policy with a verifier.  AlphaProof makes this separation
explicit: a policy and value network guide an AND--OR tree search inside Lean;
distributed actors receive adaptive compute budgets; and kernel-verified
outcomes feed back into training
\citep{hubert2025alphaproof}.  Recent systems extend the loop to iterative
conjecturing and proving \citep{dong2025stp}, natural-language research agents
\citep{feng2026aletheia}, and large-scale formal attacks on open problems
\citep{tsoukalas2026formalresearch}.

These systems make a resource question unavoidable.  Consider a research
architecture attacking many mathematical targets in parallel.  Each target
may admit several routes---sieve, geometric, spectral, probabilistic, or
computational---and several tools, including language models, computer
algebra, SAT or SMT solvers, specialized domain engines, and proof assistants.
Expert judgment may favor one route, but that judgment should change its
weight rather than silently erase every alternative.  A failed search should
lower confidence in the route that failed, not become evidence that the target
is false.  Likewise, a fast exploratory engine should remain usable without
turning its unchecked output into a mathematical premise.

The central problem is therefore not merely to generate more candidates.  It
is to allocate finite resources over targets, routes, and tools while
preserving three different properties:

\begin{enumerate}[leftmargin=*,itemsep=2pt]
  \item \emph{search coverage}: a finite valid proof should not become
        permanently unreachable because a learned policy assigns its route
        zero attention;
  \item \emph{evidentiary soundness}: a candidate should change mathematical
        status only after an adequate, independently checkable procedure has
        accepted the exact object consumed downstream;
  \item \emph{economic composability}: search should account for the cost of
        verifying, transporting, and composing its outputs, not only for the
        cost of generating them.
\end{enumerate}

The Build Thesis distinguishes structural enablement from the infrastructure
that operationalizes it \citep{li2026build}.  Proof Engine Infrastructure
supplies a typed claim graph in which only earned evidence-bearing edges derive
closure \citep{li2026proofengine}.  The claim-level verification framework
binds each claim to its object, checker, scope, residual assumptions, and
version \citep{li2026auditable}; the mathematical-engineering account adds
dependency closure, resource profiles, and downstream bottlenecks
\citep{li2026engineering}.  Taking those components as fixed interfaces, this
paper adds the active layer they leave open: how to allocate generation effort
before a candidate reaches an evidentiary gate.

The paper makes five contributions.

\begin{enumerate}[leftmargin=*,itemsep=2pt]
  \item It separates the finite proof radius of a positive instance from the
        impossibility of a computable uniform radius bound.  These elementary
        results serve as boundary conditions rather than novelty claims.
  \item It defines \emph{proof mass} and \emph{policy-relative proof distance}
        for verifier-guided generation, making the value of guidance an exact
        multiplicative change in expected search work.
  \item It proves a completeness-preservation result for an arbitrary adaptive
        guided policy mixed with a support-complete baseline, together with a
        deterministic dovetailing analogue and a finite-portfolio corollary.
  \item It defines a proof-search control plane over an earned-claim hypergraph
        and proves an exact pairwise budget-balancing theorem for homogeneous
        AND routes, connecting local search response to graph-level closure.
  \item It extends the resource objective from generation cost to the total
        cost of producing independently checkable, semantically bound, and
        composable proof edges.
\end{enumerate}

Sections~2--4 establish the formal search boundary, policy-relative distance,
and completeness reserve.  Section~5 lifts these objects to earned-claim
hypergraphs and proves the allocation theorem.  Section~6 integrates
heterogeneous tools and verification boundaries.  Sections~7--10 derive system
implications, evaluation criteria, prior-work distinctions, and scope limits.

\section{The Formal Boundary: Finite Positive Radius, No Uniform Cutoff}

We begin with the minimal formal structure needed by the later control model.
The claims in this section are classical consequences of effective proof
checking, recursive enumerability, and undecidability
\citep{church1936,turing1936,dean2025difficulty}.  Their role is to identify
exactly what a search architecture may and may not guarantee.

\begin{definition}[Machine-checkable proof system]\label{def:machine-checkable}
Let $\Sigma$ be a finite alphabet.  A machine-checkable proof system $T$ is
represented here by a total decidable relation
\[
  \Verify_T(p,\varphi)\in\{0,1\},
  \qquad p,\varphi\in\Sigma^*.
\]
The theorem set is
\[
  \Thm(T)=\{\varphi:\exists p\;\Verify_T(p,\varphi)=1\}.
\]
The proof code $p$ may encode a linear derivation, a proof tree, a DAG with
shared lemmas, a certificate, or any finite artifact accepted by the declared
checker.
\end{definition}

\begin{definition}[Proof radius]\label{def:proof-radius}
For a target $\varphi$, define its proof radius relative to $T$ and the chosen
encoding by
\[
  R_T(\varphi)=
  \min\{\,|p|:\Verify_T(p,\varphi)=1\,\},
\]
with $R_T(\varphi)=\infty$ when no such proof code exists.
\end{definition}

\begin{theorem}[Finite positive radius]
\label{thm:finite-radius}
If $\varphi\in\Thm(T)$, then $R_T(\varphi)<\infty$.  Length-lexicographic
enumeration returns a shortest proof after checking at most
\[
  \sum_{k=0}^{R_T(\varphi)}|\Sigma|^k
  =\frac{|\Sigma|^{R_T(\varphi)+1}-1}{|\Sigma|-1}
\]
candidate strings when $|\Sigma|>1$.
\end{theorem}

\begin{proof}
Provability supplies at least one finite proof code, so the nonempty set of its
lengths has a least element.  There are only finitely many strings up to that
length.  A total verifier checks them one by one, and length ordering makes the
first accepted code shortest under the declared encoding.
\end{proof}

Unlike the metaphor of an infinite library containing every finite text,
Theorem~\ref{thm:finite-radius} gives a target-conditioned procedure with a
decidable acceptance test and a finite certificate at termination.  It does
not give a feasible algorithm.  The bound is exponential in an unknown radius
and says nothing about whether the required work fits within physical
resources.

\begin{proposition}[No computable uniform proof-radius bound]
\label{prop:no-bound}
Suppose $\Thm(T)$ is undecidable.  There is no total computable function
$B:\Sigma^*\to\mathbb N$ such that
\[
  \varphi\in\Thm(T)\quad\Longrightarrow\quad R_T(\varphi)\leq B(\varphi)
\]
for every target $\varphi$.
\end{proposition}

\begin{proof}
If $B$ existed, compute $B(\varphi)$ and check every proof code of length at
most that value.  Acceptance would decide that $\varphi$ is a theorem; failure
of all candidates would decide that it is not.  This would decide
$\Thm(T)$, contradicting the hypothesis.
\end{proof}

The same argument rules out a computable bound depending only on statement
length.  Every positive instance therefore has a finite target radius, but no
general procedure announces a sufficient cutoff before search.  Proof length
also depends strongly on the chosen formal system, as Gödelian speed-up results
show \citep{godel1936}.  Proof complexity studies this system-relative
efficiency directly \citep{cookreckhow1979}.

\begin{remark}[Radius is representation-relative]
A derived lemma can be installed as a reusable named component and then cited
in one local step.  This does not add a new mathematical consequence, but it
can collapse local proof length.  A research architecture must therefore
distinguish local citation distance from the transitive cost of producing,
checking, and maintaining the cited component.  The claim graph introduced
below charges the latter rather than treating a name as free mathematical
progress.
\end{remark}

Three boundaries follow.  First, a true but $T$-unprovable target has infinite
radius relative to $T$; additional compute cannot close it.  Second, a failed
search does not reveal whether the current radius is enormous or infinite.
Third, changing the language, definitions, library, or admissible tools changes
the search geometry.  Guidance should therefore be evaluated relative to a
declared system and resource horizon, not as an intrinsic measure of
mathematical truth.

\section{Guidance as Probability Redistribution}

Exhaustive enumeration treats candidate codes uniformly.  LLM-guided theorem
proving instead uses a learned proposal distribution.  Guidance therefore
changes where search probability goes, while the verifier continues to decide
which candidates are valid.

Fix a target $\varphi$ and a protocol for one budgeted attempt.  The protocol
may run a tactic tree, call tools, or construct auxiliary lemmas.  Encode its
final finite output as a candidate $p$, treating timeouts and abandoned
attempts as rejecting outcomes.  In what follows, $\pi$ denotes this complete
proposal protocol---not the language model alone---and $\mu_\pi$ is its output
distribution.

\begin{definition}[Proof mass and policy-relative distance]\label{def:proof-mass}
The valid-proof set, proof mass, and policy-relative proof distance are
\[
\begin{aligned}
  \mathcal P_T(\varphi)
    &=\{p:\Verify_T(p,\varphi)=1\},\\
  M_\pi(\varphi)
    &=\sum_{p\in\mathcal P_T(\varphi)}\mu_\pi(p),\\
  D_\pi(\varphi)
    &=-\log M_\pi(\varphi),
\end{aligned}
\]
with $D_\pi(\varphi)=\infty$ when $M_\pi(\varphi)=0$.
\end{definition}

The definition sums over all accepted proof objects, not only a single
canonical proof.  A policy can therefore improve by concentrating on one good
route or by distributing mass across several independent routes.

\begin{proposition}[Verifier-guided hitting law]
\label{prop:hitting}
Run independent attempts from a fixed policy $\pi$ and verify every output.
If $M=M_\pi(\varphi)>0$, then the number $N$ of attempts through the first valid
proof is geometric:
\[
  \E[N]=\frac{1}{M},
  \qquad
  \Prob(N>n)=(1-M)^n\leq e^{-nM}.
\]
Consequently, $n\geq \log(1/\delta)/M$ attempts suffice for success probability
at least $1-\delta$.
\end{proposition}

\begin{proof}
Each attempt succeeds exactly when its output lies in
$\mathcal P_T(\varphi)$, an event of probability $M$.  The claims are the
standard geometric-distribution identities and the inequality
$1-M\leq e^{-M}$.
\end{proof}

\begin{proposition}[Cost-aware hitting law]
\label{prop:cost-hitting}
Let $(C_j,S_j)$ be independent, identically distributed pairs of attempt cost
and success indicator.  Assume $0<M=\Prob(S_j=1)$ and
$\E[C_j]<\infty$.  Cost may still be correlated with success within an
attempt.  If $N$ is the first successful attempt, then
\[
  \E\!\left[\sum_{j=1}^{N}C_j\right]
  =\frac{\E[C_1]}{M}.
\]
\end{proposition}

\begin{proof}
The expected total equals the expected successful-attempt cost plus the
expected number $(1-M)/M$ of failures times the expected failure cost:
\[
 \E[C\mid S=1]
 +\frac{1-M}{M}\E[C\mid S=0]
 =\frac{\E[C]}{M}.
\]
\end{proof}

Now suppose a sequential protocol generates a proof path
$\rho=(a_1,\ldots,a_L)$ through states $s_1,\ldots,s_L$.  Its path probability
and path surprisal are
\[
  \mu_\pi(\rho)=\prod_{j=1}^{L}\pi(a_j\mid s_j),
  \qquad
  d_\pi(\rho)=-\sum_{j=1}^{L}\log\pi(a_j\mid s_j).
\]
If this path dominates the target's proof mass, the expected attempt count is
approximately $e^{d_\pi(\rho)}$.  Uniform search with effective branching
factor $b$ gives $d=L\log b$, recovering the familiar $b^L$ explosion.  Proof
length alone is therefore an incomplete search distance.  A longer proof with
substantial policy mass may be easier to discover than a short proof whose
essential construction receives almost none.

\begin{corollary}[Multiplicative value of guidance]
\label{cor:guidance}
For two fixed policies with positive proof masses, define the guidance gain
from $\pi$ to $\pi'$ by
\[
  G(\pi\to\pi';\varphi)
  =D_\pi(\varphi)-D_{\pi'}(\varphi)
  =\log\frac{M_{\pi'}(\varphi)}{M_\pi(\varphi)}.
\]
Then
\[
  \frac{\E[N_\pi]}{\E[N_{\pi'}]}=e^G.
\]
Measured in bits, $G/\log 2$ is the base-two reduction in effective search
work.
\end{corollary}

This identifies a sharper model-training objective than single-step tactic
accuracy.  What matters for end-to-end discovery is the total mass retained on
at least one complete valid proof.  A locally plausible policy can still have
infinite target distance if one necessary action receives zero support.

\section{Completeness-Preserving Guided Search}

Guidance is valuable because uniform enumeration is expensive.  Guidance is
dangerous because a learned or human-weighted policy can permanently prune the
only valid route.  The remedy is not to abandon guidance, but to separate a
high-throughput guided lane from a small completeness reserve.

\begin{definition}[Support-complete baseline]\label{def:support-complete}
A family of attempt policies $\pi^{\mathrm f}_\varphi$ is support-complete for
$T$ if
\[
  \varphi\in\Thm(T)
  \quad\Longrightarrow\quad
  M_{\pi^{\mathrm f}_\varphi}(\varphi)>0.
\]
One construction samples a proof-code length from a full-support distribution,
samples uniformly at that length, and invokes $\Verify_T$.  A deterministic
alternative dovetails the length-lexicographic enumeration of
Theorem~\ref{thm:finite-radius}.
\end{definition}

Let $\pi^{\mathrm g}_{t,\varphi}$ be any adaptive guided policy at attempt $t$.
It may depend on human advice, model updates, prior failures, retrieved
literature, claim-graph state, and tool feedback.  For fixed
$\varepsilon\in(0,1]$, define
\[
  \pi_{t,\varphi}^{\varepsilon}
  =(1-\varepsilon)\pi^{\mathrm g}_{t,\varphi}
   +\varepsilon\pi^{\mathrm f}_\varphi.
\]

\begin{theorem}[Completeness preservation under arbitrary adaptive guidance]
\label{thm:fair-mixture}
Let $\varphi\in\Thm(T)$, and write
$m_{\mathrm f}=M_{\pi^{\mathrm f}_\varphi}(\varphi)>0$.  At each attempt,
select the support-complete lane independently with probability
$\varepsilon>0$; otherwise use an arbitrary history-dependent guided lane.
Then
\[
  \Prob(\text{no valid proof in the first }n\text{ attempts})
  \leq(1-\varepsilon m_{\mathrm f})^n,
\]
so a valid proof is found almost surely and
\[
  \E[N]\leq\frac{1}{\varepsilon m_{\mathrm f}}.
\]
The guided policy may change without restriction and may itself assign zero
mass to every valid proof.
\end{theorem}

\begin{proof}
Conditioned on any previous history without success, the next attempt succeeds
with probability at least $\varepsilon m_{\mathrm f}$ through the baseline
lane.  Multiplying the conditional failure bounds gives the displayed
inequality.  The tail tends to zero, and summing it gives the expectation
bound.
\end{proof}

\begin{corollary}[Distance overhead of the completeness reserve]\label{cor:reserve-distance}
For a fixed guided policy $\pi^{\mathrm g}$,
\[
  M_{(1-\varepsilon)\pi^{\mathrm g}+\varepsilon\pi^{\mathrm f}}(\varphi)
  \geq \varepsilon M_{\pi^{\mathrm f}}(\varphi),
\]
hence
\[
  D_{(1-\varepsilon)\pi^{\mathrm g}+\varepsilon\pi^{\mathrm f}}(\varphi)
  \leq D_{\pi^{\mathrm f}}(\varphi)-\log\varepsilon.
\]
The completeness insurance costs at most a factor $1/\varepsilon$ relative to
the baseline guarantee and can be much faster when guidance is useful.
\end{corollary}

The almost-sure conclusion is not a finite physical-time promise.  It assumes
that attempts continue, that the proof is expressible in the covered system,
and that the declared verifier and tools remain available.  Its significance
is architectural: adaptive weighting need not trade away asymptotic coverage.

\begin{proposition}[Deterministic dovetailing guarantee]
\label{prop:dovetail}
Let a deterministic baseline enumerator find a proof after $N_{\mathrm f}$ of
its own steps.  If a mixed scheduler allocates at least one baseline step in
every block of $q$ total steps, then the proof is found within
$qN_{\mathrm f}$ total steps, regardless of all guided work between baseline
steps.
\end{proposition}

\begin{proof}
After $qN_{\mathrm f}$ total steps the scheduler has executed at least
$N_{\mathrm f}$ baseline steps, including the accepting one.
\end{proof}

\begin{corollary}[Finite target portfolio]\label{cor:finite-portfolio}
Consider finitely many active targets $\varphi_1,\ldots,\varphi_n$.  Assign each
target a positive completeness weight $w_i$ with $\sum_iw_i=1$, and reserve a
positive total fraction $\varepsilon$ for a fair scheduler realizing those
weights.  Every $T$-provable target receives unbounded cumulative baseline work
and is eventually found.  Target weights change discovery latency, not the set
of asymptotically covered positive instances.
\end{corollary}

Together, these results formalize the role of mathematical advice.  An expert may
assign $60\%$ of guided compute to a sieve route, $25\%$ to a geometric route,
and $15\%$ to computation.  The agent may update these guided-lane weights as
evidence arrives, including setting one temporarily to zero.  Completeness is
preserved exactly when the total mixed policy retains support through its fair
lane; a guided-lane exclusion is therefore a priority decision, not a theorem
about impossibility.

\section{From a Single Proof Tree to an Earned-Claim Hypergraph}

Single-target tactic search is only one part of a research architecture.
Research portfolios also contain shared lemmas, alternative routes,
conjunctive dependencies, heterogeneous artifacts, and intermediate results
whose value lies in the downstream targets they unlock.

\begin{definition}[Earned-claim graph]\label{def:earned-graph}
At time $t$, let
\[
  \mathcal G_t=(\mathcal C,\mathcal E_t)
\]
be a typed directed hypergraph.  Each edge $e$ has a finite input set $I_e$, an
output claim $o_e$, and an evidentiary state.  Given declared roots
$R_0\subseteq\mathcal C$, only earned edges propagate reachability:
\[
  R_{k+1}=R_k\cup
  \{o_e:e\in\mathcal E_t^{\mathrm{earned}},\ I_e\subseteq R_k\},
  \qquad
  \Reach(\mathcal G_t)=\bigcup_{k\geq0}R_k.
\]
Multiple incoming edges to one output have OR semantics; multiple inputs of
one edge have AND semantics.
\end{definition}

This is the minimal portion of Proof Engine Infrastructure needed here
\citep{li2026proofengine}.  A generated proof sketch, a solver log, or even a
locally accepted artifact remains nonpropagating until the exact output claim,
assumptions, coverage, and verification boundary satisfy the declared
admissibility rule.  Search beliefs can move continuously; earned reachability
moves only through checked transitions.

\begin{definition}[Search-control action]\label{def:control-action}
A control action is a tuple
\[
  \alpha=(i,e,r,m,b),
\]
Here $i$ identifies a target in the active portfolio; $e$ is an open claim
edge; $r$ is a mathematical route; $m$ is a generation or analysis tool; and
$b$ is the planned verification boundary.  The action may yield no progress,
edge-relative negative evidence, a candidate artifact, or an artifact that
passes the checks required to earn $e$.
\end{definition}

Provability alone is too coarse a scheduling prior.  A claim may be highly
likely to be provable and still remain unreachable within any relevant budget.
Instead, define the budgeted earned-closure distribution
\[
  F_{t,\alpha}(B)
  =\Prob(
     e\text{ becomes earned within additional budget }B
     \mid \mathcal H_t,\mathcal G_t,\alpha),
\]
where $\mathcal H_t$ records human route assessments, previous attempts, tool
behavior, and model state.  This distribution combines uncertainty about the
route, proof radius, tool suitability, and verification cost.  It guides the
scheduler; it never replaces evidence.

Let $U(S)$ be a declared utility on reachable claim sets.  A simple terminal
utility is
\[
  U(S)=\sum_i u_i\ind\{T_i\in S\},
\]
but intermediate claims may receive value for reducing a minimal open cut,
serving several targets, or creating a reusable theorem component.  The
realized reachability gain of an earned edge is
\[
  \Delta U_t(e)
  =U(\Reach(\mathcal G_t\cup\{e\}))-U(\Reach(\mathcal G_t)).
\]
Before downstream premises are ready this immediate gain may be zero, so a
practical scheduler also estimates potential gain through remaining cuts.

\begin{principle}[Closure-aware resource allocation]\label{prn:closure-aware}
Allocate guided resources using the probability that an action produces an
earned edge within budget, the edge's expected downstream reachability value,
and the total cost of making the output independently consumable.  Do not rank
actions by candidate-generation rate alone.
\end{principle}

For an additional budget interval $\Delta B$, a heuristic index consistent
with this principle is
\[
  I_t(\alpha;\Delta B)
  =
  \frac{
    \widehat{\Delta U}_t(e)F_{t,\alpha}(\Delta B)
  }{
    \E[C_{\mathrm{gen}}+C_{\mathrm{verify}}+C_{\mathrm{transport}}
       +C_{\mathrm{compose}}\mid\alpha]
  }
  +\beta_t\,\mathsf{Explore}_t(\alpha).
\]
The exploration term expresses epistemic value and prevents premature route
collapse inside the guided lane; the separate completeness reserve protects
formal coverage even when the learned index is wrong.

The portfolio-level objective over a total budget $B$ is
\[
  \max_{\Pi}
  \E_\Pi\!\left[U(\Reach(\mathcal G_B))\right]
  -\lambda\E_\Pi[C_{\mathrm{total}}],
\]
subject to the completeness allocation and to the rule that only earned edges
enter $\Reach$.  This is deliberately an objective, not a claim that all terms
already admit reliable universal calibration.

The graph exposes allocation effects that a list of independent theorem
prompts cannot represent.  A difficult lemma on many target cuts may deserve
more compute than an easy isolated theorem.  Two agents exploring distinct
incoming edges may be valuable; two agents unknowingly repeating one route may
not be.  A failure lowers the posterior of that route and records its
obstruction, while other incoming edges to the same claim remain open.

\subsection{An exact allocation law for a homogeneous AND route}
\label{sec:and-allocation}

The general objective does not itself determine an allocation rule.  The
following tractable case does.  Consider two required edges of an AND route.
Each independent attempt on either edge fails with the same probability
$q\in[0,1]$.  After $n$ attempts, the edge has been earned with probability
$1-q^n$.  Independence is explicit here: the model applies to fresh attempts
without shared failure causes, not to proof search in general.

\begin{definition}[Homogeneous edge and two-edge route response]\label{def:and-response}
For $q\in\mathbb R$ and $n\in\mathbb N$, define
\[
  S_q(n)=1-q^n.
\]
For a nonnegative multiplier $R$ representing the already-funded remainder of
the route, define
\[
  A_q(R;a,b)=R\,S_q(a)S_q(b).
\]
When $R$ is the product of the success probabilities of the other required
edges, $A_q$ is the closure probability of the whole AND route.
\end{definition}

\begin{theorem}[One-unit balancing identity]\label{thm:and-balance-identity}
For every $q,R\in\mathbb R$ and $a,d\in\mathbb N$,
\[
\begin{aligned}
 &A_q(R;a+1,a+d+1)-A_q(R;a,a+d+2)\\
 &\qquad =R\,q^a(1-q)(1-q^{d+1}).
\end{aligned}
\]
\end{theorem}

\begin{proof}
Expand both products.  The $q^{2a+d+2}$ terms cancel, and the two remaining
differences factor as
\[
R(q^a-q^{a+1}-q^{a+d+1}+q^{a+d+2})
=R\,q^a(1-q)(1-q^{d+1}).
\]
\end{proof}

\begin{theorem}[Strict improvement of an underfunded conjunct]\label{thm:and-balance-strict}
If $R>0$, $0<q<1$, and one required edge has at least two more attempts than
the other, transferring one attempt from the better-funded edge to the
underfunded edge strictly increases AND-route closure probability while
preserving total search budget.  Concretely, for all $a,d\in\mathbb N$,
\[
  A_q(R;a,a+d+2)<A_q(R;a+1,a+d+1).
\]
\end{theorem}

\begin{proof}
Theorem~\ref{thm:and-balance-identity} gives the gain.  Each factor on its
right-hand side is strictly positive: $R>0$, $q^a>0$, $1-q>0$, and
$1-q^{d+1}>0$ because $d+1>0$ and $0<q<1$.
\end{proof}

\begin{corollary}[Pairwise balance of an optimal homogeneous AND allocation]
\label{cor:and-pairwise-balance}
Fix $R>0$ and $0<q<1$.  If $(a,b)$ maximizes $A_q(R;a,b)$ among all pairs of
natural-number budgets with the same total $a+b$, then
\[
  a\leq b+1\qquad\text{and}\qquad b\leq a+1.
\]
Equivalently, an optimal pair differs by at most one attempt.
\end{corollary}

\begin{proof}
If, for example, $b\geq a+2$, write $b=a+d+2$.  Theorem
\ref{thm:and-balance-strict} constructs the same-total allocation
$(a+1,b-1)$ with strictly larger closure probability, contradicting
optimality.  The other inequality is symmetric.
\end{proof}

The theorem is local, but its use is not limited to two-edge problems.  In a
larger homogeneous AND route, $R$ is the product contributed by all other
positively funded required edges.  Once those edges have positive success
probability, no pairwise-optimal scheduler can give one conjunct at least two
more attempts than another.  Heterogeneous hazards, shared failure modes, OR
routes, and reusable cross-target producers require the fuller control model.

\section{Heterogeneous Tools and Verification Boundaries}

The control plane must not equate proof search with a single proof assistant.
A high-frequency exploratory loop may use a domain-specific geometry engine,
a rewriting system, a CAS, a SAT solver, executable search, or an LLM.  Lean,
Isabelle, Coq, a small certificate checker, or direct witness validation may
serve at different cut points.  Tool uniformity is not the invariant.  The
checked object must instead be adequate for the exact claim and composable with
its consumers.

For a claim $C$ and verification boundary $b$, write $V_b(C)$ for the cost of
independent checking.  For an edge $e$, write $K(e)$ for the cost of
establishing that its checked inputs jointly support its exact output.  A
target-supporting earned subgraph $\mathcal G_T$ has verification-composition
cost
\[
  \sum_{C\in\mathcal C(\mathcal G_T)}V_{b(C)}(C)
  +\sum_{e\in\mathcal E(\mathcal G_T)}K(e),
\]
following the graph-level objective in Proof Engine Infrastructure
\citep{li2026proofengine}.  The present control plane adds the upstream cost of
discovering those components and treats boundary choice as part of each search
action.

\begin{principle}[Proof-carrying exploration]\label{prn:proof-carrying}
Use the fastest suitable engine in the exploratory loop only when its positive
outputs can be converted into an independently checkable artifact whose
semantics match the claim consumed downstream.
\end{principle}

For example, a specialized geometry engine may derive millions of relations
and return a compact construction trace.  A narrow checker can validate that
trace.  Depending on the required assurance, a proof assistant may then verify
the checker, the trace, selected semantic interfaces, or the final theorem.
There is no need to elaborate every failed branch in a heavyweight
environment.  Conversely, an expensive search does not make the fast engine's
unchecked assertion earned.

The feedback loop is therefore:
\[
\begin{array}{c}
\text{human targets, utilities, and route priors}\\
\downarrow\\
\text{target--edge--route--tool scheduler}\\
\downarrow\\
\text{parallel candidate generation and analysis}\\
\downarrow\\
\text{claim-relative verification profiles}\\
\downarrow\\
\text{candidate-to-earned gate and claim-graph closure}\\
\downarrow\\
\text{updated blockers, failures, shared lemmas, and resource posteriors}\\
\circlearrowleft
\end{array}
\]

The two planes answer different questions.  The control plane asks where to
search next; the evidentiary plane asks what the resulting artifact
establishes.  No probability, model identity, agent reputation, or quantity of
compute can substitute for the latter.

\section{Implications for LLM Mathematics}

\subsection{Train for complete-path mass, not only local accuracy}

Tactic accuracy, pass@$k$, and benchmark solve rate remain useful, but they do
not explain why a target was missed.  A control-plane evaluation should track
probability mass along complete successful paths, identify necessary
low-probability actions that were pruned, and measure verified work per unit of
total cost.  Corollary~\ref{cor:guidance} predicts an exponential payoff from
modest reductions in cumulative path surprisal.

This interpretation is consistent with current policy/value proof search
\citep{hubert2025alphaproof,yang2023leandojo}: the network is valuable not
because it certifies a tactic, but because it concentrates finite search on
productive branches.  The distinction becomes more important at the research
frontier, where the relevant path may contain method changes, literature
retrieval, new definitions, auxiliary computation, and specialized tools
rather than only tactics from an established library.

\subsection{Calibrate closure within budget, not timeless confidence}

A useful model should estimate a survival curve or hazard for earned closure.
That estimate asks how likely an edge is to close in the next unit of compute,
given that it has not closed so far.  The model should represent proof,
refutation, formalization mismatch, and unresolved status separately.  In an
undecidable setting, a timeout cannot convert ``no proof found'' into a final
unprovability label.

Human guidance enters as a prior over routes and semantic interpretations.
Machine experience updates it.  Exact refutations, counterexamples, and
independence results may justify hard exclusions; ordinary search failure
usually justifies a weight change.  This is a probabilistic control rule, not a
probabilistic definition of mathematical truth.

\subsection{Coordinate portfolios through shared cuts}

Parallelism is most valuable when agents attack distinct bottlenecks whose
outputs compose.  The scheduler should expose ownership, route identity,
consumer claims, and verification plans before launching work.  Shared
producer lemmas receive portfolio value; semantically duplicate fragments do
not receive artificial credit merely because several agents generated them.
Independent replication remains valuable when resources are deliberately
allocated to verification rather than accidentally duplicated in discovery.

\subsection{Do not optimize generation into a downstream queue}

AI-assisted mathematics can make candidate production grow faster than
dependency closure, replay, exposition, and maintenance.  The resulting
bottleneck displacement is developed in the mathematical-engineering account
\citep{li2026engineering}.  The control plane incorporates that lesson by
pricing verification and composition before allocating search.  A route that
produces a huge opaque artifact may be mathematically valid yet economically
inferior to a route that produces a reusable checked component.

\section{Evaluation Program}

The framework yields measurements that differ from model-only leaderboards.
For each target portfolio and resource horizon, report:

\begin{enumerate}[leftmargin=*,itemsep=2pt]
  \item earned target closure per GPU-hour, dollar, and verifier-hour;
  \item calibration of $F_{t,\alpha}(B)$ against realized earned closure;
  \item proof mass or a controlled estimator of complete-path mass, together
        with realized search work;
  \item route diversity and the frequency of necessary actions pruned by the
        guided lane;
  \item reachability gain and reuse count of newly earned edges;
  \item generation, verification, transport, and composition cost separately;
  \item accumulated verification debt: candidates awaiting an adequate check,
        semantic binding, or dependency completion;
  \item the fraction of closures found through the guided lane and through the
        completeness reserve.
\end{enumerate}

The following three hypotheses make the evaluation program testable.

\begin{conjecture}[Guidance-gain calibration]\label{conj:guidance-calibration}
Holding verifier and attempt protocol fixed, reductions in estimated
policy-relative proof distance should predict multiplicative reductions in
realized candidate count according to Corollary~\ref{cor:guidance}, up to
dependence, truncation, and cost heterogeneity that must be reported.
\end{conjecture}

\begin{conjecture}[Graph-aware portfolio gain]\label{conj:graph-portfolio}
At equal total compute, scheduling by expected earned reachability gain should
close more weighted portfolio targets than scheduling each theorem
independently, especially when targets share load-bearing producers.
\end{conjecture}

\begin{conjecture}[Engineering complementarity]\label{conj:engineering-complementarity}
The marginal return from stronger LLM guidance should be larger when candidate
outputs have predeclared verification boundaries, semantic interfaces, and
claim-graph consumers than when proof generation is optimized without those
complements.
\end{conjecture}

These conjectures are falsifiable.  Proof-mass estimates may prove too
unstable; graph maintenance may cost more than it saves; or specialized tools
may create transport debt that dominates their search advantage.  The
framework requires those costs to be measured rather than assumed away.

\section{Relation to Prior Work}

\paragraph{Computability and proof length.}
Recursive enumerability and the negative solution to the general decision
problem supply the classical positive-instance/unknown-cutoff boundary
\citep{church1936,turing1936}.  Gödelian speed-up and proof complexity show
that proof length depends strongly on the formal system
\citep{godel1936,cookreckhow1979}.  Dean and Naibo apply computability and
complexity limits directly to contemporary AI and mathematical discovery
\citep{dean2025difficulty}.  Section~2 uses these results as operational
boundary conditions; it does not claim them as new.

\paragraph{Universal and heuristic search.}
Levin search allocates resources across candidate programs with a universal
length prior.  Hutter combines program execution with enumeration of proofs of
correctness and runtime bounds \citep{levin1973,hutter2002}.  Those systems
give broad guarantees for well-defined search problems.  The present framework
is narrower in universality but broader in research organization: it combines
adaptive human/LLM route weights, heterogeneous verification boundaries,
multiple mathematical targets, and downstream claim composition.
Policy-relative distance serves here as an operational accounting quantity,
not as a replacement for algorithmic probability or Kolmogorov complexity.

\paragraph{Portfolio selection and strategy scheduling.}
Classical ATP systems have long allocated time across heterogeneous strategies
rather than relying on one universal configuration.  MaLARea couples learning
and reasoning over large theorem libraries \citep{urban2014malarea}.  Modern
Vampire schedules discover and regularize complementary strategies for
sequential or parallel execution \citep{bartek2024scheduling}.  More generally,
portfolio-based algorithm selection has learning-theoretic guarantees and
overfitting limits \citep{balcan2021portfolio}.  These systems directly
precede weighted tool and strategy choice.  Our control object adds three
joint requirements: allocation ranges over typed mathematical claims rather
than only prover configurations; success means an earned, semantically bound
edge rather than a solver return alone; and utility includes downstream
reachability together with verification and composition costs.

\paragraph{Automated theory formation.}
AM, HR, MATHsAiD, and related theory-exploration systems generate concepts,
conjectures, proofs, and counterexamples rather than solving only a supplied
theorem \citep{lenat1976,colton2002,mccasland2017}.  Recent self-play theorem
provers continue this direction \citep{dong2025stp}; Chen and Li model
theorems as a semantic-similarity graph and study conditions for exponential
self-play growth \citep{chenli2026selfplay}.  Their graph supports curriculum
expansion.  The graph here instead records typed evidentiary dependency and
earned reachability; its control problem is which claim edge to fund and how
to verify it.

\paragraph{Neural and agentic theorem proving.}
Generative proof models, retrieval-augmented provers, and large-scale formal
RL systems learn proposal policies and guide tree search
\citep{polu2020,yang2023leandojo,hubert2025alphaproof}.  AlphaProof already
contains policy/value guidance, progressive sampling, AND nodes, a distributed
matchmaker, and adaptive compute budgets.  Aletheia iterates generation,
verification, and revision in natural language \citep{feng2026aletheia}; a
recent formal study attacks hundreds of open problems with LLM--Lean agents
and reports per-problem costs \citep{tsoukalas2026formalresearch}.  These works
establish the inner-search precedent.  The present contribution is the outer,
completeness-preserving and verification-cost-aware control plane over a
portfolio of exact claim edges.

\paragraph{Proof Engine and mathematical engineering.}
The companion architecture separates generation from evidence and derives
closure through earned claim edges \citep{li2026proofengine}; the claim-level
framework defines heterogeneous verification profiles
\citep{li2026auditable}; and the engineering-layer paper describes dependency
closure, component contracts, resource profiles, and downstream bottleneck
movement \citep{li2026engineering}.  This paper consumes only their minimal
interfaces.  It supplies the upstream scheduler and closes the feedback loop
from graph state to future search allocation.

\section{Limitations and Scope}

The results do not imply that every true mathematical statement will be found.
They apply to targets with finite proof objects in the covered formal system.
They do not select sound new axioms, settle semantic truth beyond the system,
or give a finite physical budget for an unknown proof radius.

Policy-relative distance refers to the complete proposal protocol: model
policy, search algorithm, decoding parameters, representation, tool set,
verifier, and resource cutoff.  Changing any of these changes $\mu_\pi$; the
notation does not assert model-only invariance.  Exact proof mass is generally
unavailable and must be estimated.  Correlation between parallel attempts,
tree reuse, online learning, and variable verification costs can also make
independent-sampling predictions inaccurate.  Proposition~\ref{prop:hitting}
is therefore a baseline law, not a complete performance model.

The scheduler objective contains human choices.  Target utilities, acceptable
trust boundaries, semantic identity, and the value of explanation are not
deduced from proof syntax.  Human judgment may itself be wrong, but making its
role explicit is preferable to hiding it inside prompts or model scores.

Finally, the completeness reserve is an asymptotic insurance policy.  A tiny
reserve can be practically irrelevant within a finite project, while a large
reserve can consume resources better spent on informed search.  Choosing
$\varepsilon$ is a resource-allocation decision whose empirical optimum will
vary by domain.  The theorem guarantees what is preserved, not what fraction
is economically optimal.

\section{Conclusion}

A target provable in a declared machine-checkable system has a finite proof
radius.  For an undecidable theorem set, no computable universal cutoff reveals
that radius in advance.  AI-assisted mathematics therefore faces finite
positive witnesses, potentially enormous and unknowable discovery cost, and no
sound timeout-based test for universal failure.

LLMs, human mathematical guidance, and specialized tools address this problem
by redistributing probability and compute.  Their success is measured by the
mass they place on complete valid proof objects and by the cost of turning
those objects into earned, composable claims.  A completeness reserve allows
arbitrary adaptive guidance without permanently deleting finite proofs from
the search support.  An earned-claim hypergraph then converts local results
into a portfolio-level control problem: fund the target--route--tool actions
most likely to create valuable verified reachability, while preserving a
small fair lane and pricing the downstream verification burden.

This division of labor defines the proposed research method.  Humans specify
targets, meanings, utilities, and informed route priors.  Agents navigate and
update the weighted proof space.  Specialized engines search at the speed
appropriate to their objects, while independent checkers determine evidentiary
status.  The claim graph records what has actually become reachable.  More
compute enlarges the practical frontier; guidance changes how quickly it is
reached; verification prevents speed from becoming authority.

\section*{Declaration of Generative AI Use}

Generative AI tools assisted with literature discovery, structural drafting,
and copy editing.  The author selected the framework, mathematical statements,
scope restrictions, and final text, and takes responsibility for the paper.

\bibliographystyle{plainnat}
\bibliography{references}

\end{document}
